\documentclass[10pt]{article}
\usepackage[T1]{fontenc}
\usepackage[utf8]{inputenc}
\usepackage[french]{babel}
\usepackage{amsmath,amssymb,mathrsfs,mathpazo}
\usepackage{geometry}
\geometry{a4paper,margin=2cm}
\usepackage{graphicx,array,tabularx,tabularray,colortbl,multirow,multicol,booktabs,enumitem,pifont,ulem,fancybox,diagbox,cancel}
\usepackage{pstricks,pstricks-add,pst-plot,pst-tree,pst-node,pst-text,pst-func,pst-math,pst-bspline,pst-3dplot,pst-eucl,pst-solides3d}
\let\pscalculate\undefined
\usepackage{pst-bezier}
\usepackage{tikz,pgfplots,tkz-tab}
\usetikzlibrary{arrows,arrows.meta,patterns,calc,positioning,trees,decorations.markings,shapes.geometric,shapes.misc,shapes.arrows,intersections,decorations.pathmorphing,decorations.pathreplacing,backgrounds,shadows,babel}
\pgfplotsset{compat=1.18}
\usepgfplotslibrary{fillbetween}
\usepackage[np]{numprint}
\usepackage{listings,makecell,pas-tableur,twoopt}
\let\barre\undefined
\usepackage[pstricks]{tablvar}
\definecolor{darkgreen}{RGB}{0,100,0}
\providecommand{\red}{\color{red}}
\providecommand{\blue}{\color{blue}}
\providecommand{\green}{\color{green}}
\providecommand{\black}{\color{black}}
\providecommand{\white}{\color{white}}
\providecommand{\gray}{\color{gray}}
\providecommand{\cyan}{\color{cyan}}
\providecommand{\magenta}{\color{magenta}}
\providecommand{\yellow}{\color{yellow}}
\providecommand{\violet}{\color{violet}}
\providecommand{\orange}{\color{orange}}
\providecommand{\euro}{\texteuro}
\providecommand{\EUR}{\texteuro}
\providecommand{\ieme}{\textsuperscript{e}}
\providecommand{\degres}{\ensuremath{^\circ}}
\providecommand{\degree}{\ensuremath{^\circ}}
\providecommand{\vv}[1]{\overrightarrow{#1}}
\providecommand{\Coord}[1]{#1}
\providecommand{\decosix}{}
\providecommand{\starredbullet}{\ensuremath{\bullet}}
\pagestyle{empty}
\setlength{\parindent}{0pt}
\usepackage{adjustbox}\DeclareUnicodeCharacter{2236}{:}\DeclareUnicodeCharacter{2190}{\ensuremath{\leftarrow}}
\providecommand{\R}{\mathbb{R}}
\providecommand{\N}{\mathbb{N}}
\providecommand{\Z}{\mathbb{Z}}
\providecommand{\Q}{\mathbb{Q}}
\providecommand{\D}{\mathbb{D}}
\providecommand{\C}{\mathbb{C}}
\providecommand{\e}{{\mathrm{e}}}
\providecommand{\vect}[1]{\overrightarrow{#1}}
\providecommand{\vectt}[1]{\overrightarrow{\mathrm{#1}}}
\providecommand{\vvt}[1]{\overrightarrow{\mathrm{#1}}}
\providecommand{\np}[2][]{{#2~#1}}
\providecommand{\vv}[1]{{\overrightarrow{#1}}}
\renewcommand{\d}{{\mathrm{d}}}
\providecommand{\degres}{{^\circ}}
\providecommand{\degre}{{^\circ}}
\providecommand{\blue}{\color{blue}}
\providecommand{\red}{\color{red}}
\providecommand{\enskip}{\quad}
\providecommand{\llbracket}{\lbrack\!\lbrack}
\providecommand{\rrbracket}{\rbrack\!\rbrack}
\providecommand{\euro}{{\text{€}}}
\providecommand{\no}{{\text{n°}}}
\providecommand{\interval}[2]{{[#1;#2]}}
\providecommand{\intervFO}[2]{{[#1;#2[}}
\providecommand{\dotfill}{\ldots}

\providecommand{\ds}{\displaystyle}
\providecommand{\textding}[1]{#1}
\providecommand{\Oij}{\(\mathrm{O};\vec{\imath},\vec{\jmath}\)}
\providecommand{\Oijk}{\(\mathrm{O};\vec{\imath},\vec{\jmath},\vec{k}\)}
\providecommand{\Ouv}{\(\mathrm{O};\vec{u},\vec{v}\)}
\providecommand{\intervFF}[2]{\ensuremath{[#1;#2]}}
\renewcommand{\euro}{\ifmmode\text{\texteuro}\else\texteuro\fi}
\DeclareUnicodeCharacter{00B0}{\ifmmode^{\circ}\else\textdegree\fi}
\DeclareUnicodeCharacter{20AC}{\ifmmode\text{\texteuro}\else\texteuro\fi}
% Énoncés : Éducation nationale. Transcription et corrigés : APMEP, Denis Vergès et contributeurs. Découpage et classement : CoopMaths, Sébastien Lozano, Éric Elter et contributeurs. Adaptation de la présentation : Maths Explorer.
% CC BY-SA : https://coopmaths.fr/www/about/
% https://coopmaths.fr/alea/static/eam/2026/tex/eam_specialite_2026_06_metropole_3_cor.tex
\begin{document}
\noindent\textit{APMEP / CoopMaths. CC BY-SA. Présentation adaptée par Maths Explorer.}\par\medskip
Dans le plan muni d'un repère orthonormal, on considère les points \mbox{P~$(4;0)$} et \mbox{K~$(1;0)$}.

On considère le réel $x$ et on note M le point de coordonnées \mbox{$(x;3)$.}
\begin{center}

	\begin{tikzpicture}[xscale=1,yscale=1,rotate=0,every node/.style={scale=1,rotate=0}]
	\clip(-0.6,-0.6)rectangle(4.6,3.6);

	\draw[blue!30,thick](0,0)grid(4.6,3.6);
	\draw[blue, line width=0.3mm,->, >=triangle 45](-0.6,0)--(4.6,0);
	\draw[blue, line width=0.3mm,->, >=triangle 45](0,-0.6)--(0,3.6);
	\foreach \i in {1,2,...,4}{\draw(\i,-0.2)node{\textbf{\i}};}
		\foreach \j in {1,2,3}{\draw(-0.2,\j)node{\textbf{\j}};}
		\draw(-0.2,-0.2)node{\textbf{0}};
	\draw[fill=black](1,0)circle(0.04)node[above right]{\textbf{K}}(4,0)circle(0.04)node[above right]{\textbf{P}}(2.3,3)circle(0.04)node[above]{\textbf{M}};

\end{tikzpicture}

\end{center}
\begin{enumerate}
	\item Donner les coordonnées du vecteur \mbox{$\vect{\text{KP}}$} ainsi que sa norme.


	\textcolor{red}{\textsc{Correction}}


	Les coordonnées de $\vect{\text{KP}}$ sont 	\mbox{$(3;0)$}
	 et sa norme est $3$.
	\item Exprimer en fonction de $x$ les coordonnées du vecteur \mbox{$\vect{\text{KM}}$} ainsi que sa norme.




	\textcolor{red}{\textsc{Correction}}


	Les coordonnées de $\vect{\text{KM}}$ sont \mbox{$(x-1;3)$}

	et sa norme est \quad \mbox{$\sqrt{(x-1)^2+3^2}=\sqrt{(x-1)^2+9}=\sqrt{x^2-2x+10}$}.

	\item Montrer que le produit scalaire \mbox{$\vect{\text{KP}}\cdot\vect{\text{KM}}$} est égal à \mbox{$3x-3$.}

		\textcolor{red}{\textsc{Correction}}

		On connait les coordonnées des deux vecteurs : \quad \mbox{$\vect{\text{KP}}\cdot\vect{\text{KM}}=3(x-1)+0\times3=3x-3$}
	\item Montrer que si l'angle \ensuremath{\widehat{\text{PKM}}} est égal à \mbox{$\dfrac\pi 3$}, alors le réel $x$ est solution de l'équation

	\begin{center}
		(\textit{E})\hspace{10mm}\mbox{$\sqrt{(x-1)^2+9}=2x-2$}
	\end{center}



	\textcolor{red}{\textsc{Correction}}

	On sait aussi que\quad  \mbox{$\vect{\text{KP}}\cdot\vect{\text{KM}}=\left\|\overrightarrow{\text{KP}}\right\|\times\left\|\overrightarrow{\text{KM}}\right\|\times \cos(\widehat{\text{PKM}})$}.

	Puisque \quad \mbox{$\cos\left(\dfrac\pi 3\right)=\dfrac12$}, \quad on peut donc écrire que \quad \mbox{$\vect{\text{KP}}\cdot\vect{\text{KM}}=\dfrac12\times3\times\sqrt{(x-1)^2+9}$.}

	On a donc \quad \mbox{$\dfrac32 \sqrt{(x-1)^2+9}=3x-3\iff \sqrt{(x-1)^2+9}=2x-2$}

	lorsque l'angle \ensuremath{\widehat{\text{PKM}}} est \mbox{$\dfrac\pi 3$}, $x$ est une solution de l'équation (\textit{E}).
	\item Vérifier que le réel \mbox{$1+\sqrt3$} est solution de l'équation (\textit{E}).

		\textcolor{red}{\textsc{Correction}}

		Si $x$ est une solution de l'équation (\textit{E}), alors \quad \mbox{$(x-1)^2+9=(2x-2)^2$},

		ce qui revient à \quad \mbox{$3x^2-6x-6=0$}.

		On calcule \quad \mbox{$\Delta=36+72=108=3\times6^2$}.

		L'équation a deux solutions qui sont \quad  \mbox{$\dfrac{6+6\sqrt3}{6}=1+\sqrt3$} \quad et \quad \mbox{$1-\sqrt3$}.

		\mbox{$1+\sqrt3$} est une solution de l'équation (\textit{E}).
\end{enumerate}

\end{document}
